% 数学ノート集・共通テンプレート v2（積分論2026年度第1部の版面を参考）
% A5判。章内で定義・定理・例・注意等の番号を共通化。
% コンパイル： platex -interaction=nonstopmode math_note_template_v2.tex (2回)
%              dvipdfmx math_note_template_v2.dvi
\documentclass[a5paper,10pt,onecolumn]{jsarticle}
\usepackage[dvipdfmx,top=22mm,bottom=20mm,inner=19mm,outer=18mm,headheight=13pt,headsep=6mm]{geometry}
\usepackage{amsmath,amssymb,amsthm,mathtools}
\usepackage{enumitem,fancyhdr,titlesec}
\usepackage[dvipdfmx,bookmarks=true,colorlinks=false]{hyperref}
\usepackage{pxjahyper}

% ----- 差し替える文書情報 -----
\newcommand{\NoteTitle}{素数定理とゼータ関数の解析}
\newcommand{\NoteSeries}{数学ノート集}

% ----- 文字・組版 -----
% 参考PDFの質素な明朝＋Computer Modernの印象を引き継ぎながら，
% A5判で窮屈にならないよう本文の行間を少し確保する．
\linespread{1.13}
\setlength{\parindent}{1zw}
\setlength{\parskip}{0pt}
\setlength{\abovedisplayskip}{6pt plus 2pt minus 2pt}
\setlength{\belowdisplayskip}{6pt plus 2pt minus 2pt}
\setlength{\abovedisplayshortskip}{3pt plus 2pt minus 1pt}
\setlength{\belowdisplayshortskip}{4pt plus 2pt minus 1pt}
\setlist[enumerate]{topsep=3pt,itemsep=1pt,parsep=0pt,leftmargin=2.4zw}
\titleformat{\section}{\normalfont\normalsize\bfseries}{\thesection}{0.8zw}{}
\titlespacing*{\section}{0pt}{2.0ex plus .4ex}{.9ex}
\titleformat{\subsection}{\normalfont\normalsize\bfseries}{\thesubsection}{0.7zw}{}
\titlespacing*{\subsection}{0pt}{1.5ex plus .3ex}{.6ex}
\raggedbottom

% ----- 柱とページ番号：ページ上部、罫線なし -----
\fancypagestyle{mathnote}{%
  \fancyhf{}%
  \fancyhead[L]{\footnotesize \NoteSeries\ （\NoteTitle）}%
  \fancyhead[R]{\footnotesize \thepage}%
  \renewcommand{\headrulewidth}{0pt}%
  \renewcommand{\footrulewidth}{0pt}%
}
\pagestyle{mathnote}

% ----- 節内通し番号の定義・定理群 -----
\newtheoremstyle{mathnotetheorem}
 {4.5pt}{4.5pt}{\normalfont}{}%
 {\bfseries}{}{0.65zw}%
 {\thmname{#1}\thmnumber{ #2}\thmnote{（#3）}}
\theoremstyle{mathnotetheorem}
\newtheorem{theorem}{定理}[section]
\newtheorem{definition}[theorem]{定義}
\newtheorem{lemma}[theorem]{補題}
\newtheorem{proposition}[theorem]{命題}
\newtheorem{corollary}[theorem]{系}
\newtheorem{example}[theorem]{例}
\newtheorem{remark}[theorem]{注意}
\renewcommand{\proofname}{証明}
\renewcommand{\qedsymbol}{$\square$}
\makeatletter
\renewenvironment{proof}[1][\proofname]{%
  \par\pushQED{\qed}\normalfont\topsep4\p@\@plus2\p@\relax
  \trivlist\item[\hskip\labelsep\bfseries #1]\ignorespaces
}{\popQED\endtrivlist\@endpefalse}
\makeatother

% ----- 使うノートで適宜拡張する記号 -----
\newcommand{\R}{\mathbb{R}}
\newcommand{\norm}[1]{\lVert #1\rVert}
\newcommand{\abs}[1]{\lvert #1\rvert}

\begin{document}
\section*{読了後にできること}
本稿を読むための前提は，Chebyshev関数，部分和公式，素数定理．
本稿の具体的な到達課題は，素数定理から素数の重み付き総和の漸近を導く．

\section{素数計数とChebyshev関数}

本稿では複素解析のCauchy積分公式・留数定理，
Lebesgueの収束定理，Dirichlet級数の絶対収束を前提とする．
素数定理を，境界線上の零点排除と
Newman型のTauber定理を通して証明する．

\begin{definition}
$x\ge1$ に対し
\[
 \pi(x)=\#\{p\le x:p\text{ は素数}\},\qquad
 \vartheta(x)=\sum_{p\le x}\log p
\]
と置く．von Mangoldt関数を
\[
 \Lambda(n)=
 \begin{cases}
 \log p&(n=p^k,\ p\text{ 素数},\ k\ge1),\\
 0&(\text{その他})
 \end{cases}
\]
と定め，
\[
 \psi(x)=\sum_{n\le x}\Lambda(n)
 \]
を第二Chebyshev関数という．
\end{definition}

\begin{lemma}[素数冪による差]\label{lem:prime-powers}
$x\ge2$ に対し
\[
 0\le\psi(x)-\vartheta(x)
 \le\sum_{2\le k\le\log_2x}\vartheta(x^{1/k}).
\]
\end{lemma}
\begin{proof}
$\Lambda$ は素数冪でのみ非零なので
\[
 \psi(x)=\sum_{\substack{p^k\le x\\k\ge1}}\log p
 =\sum_{k\ge1}\sum_{p\le x^{1/k}}\log p
 =\sum_{k\ge1}\vartheta(x^{1/k}).
\]
$k=1$ の項は $\vartheta(x)$ であり，
$p^k\ge2^k$ なので和に寄与する $k$ は
$k\le\log_2x$ に限られる．
従って差の表示と不等式を得る．
\end{proof}

\begin{lemma}[Chebyshevの上界]\label{lem:chebyshev}
ある定数 $C>0$ が存在して
\[
 \vartheta(x)\le Cx,\qquad
 \psi(x)\le Cx
 \quad(x\ge2)
\]
が成り立つ．さらに
$\psi(x)-\vartheta(x)=O(\sqrt x\log x)=o(x)$ である．
\end{lemma}
\begin{proof}
$n$ を正の整数とする．
$n<p\le2n$ の各素数は二項係数
$\binom{2n}n$ を割る．
実際 $p>n$ なので $p$ は分母の $n!^2$ を割らず，
分子の $(2n)!$ をちょうど一度割る．
従って
\[
 \prod_{n<p\le2n}p\le\binom{2n}n\le2^{2n}.
\]
対数を取って
$\vartheta(2n)-\vartheta(n)\le2n\log2$ を得る．
$2^j$ ごとの区間でこれを適用して加えると
$\vartheta(2^m)\le C_0\,2^m$ となる．
任意の $x\ge2$ に対し
$2^{m-1}<x\le2^m$ を取れば
$\vartheta(x)\le\vartheta(2^m)\le2C_0x$ なので
$\vartheta(x)=O(x)$ である．

補題\ref{lem:prime-powers}とこの評価から
\[
 \psi(x)-\vartheta(x)
 \le C\sum_{2\le k\le\log_2x}x^{1/k}
 \le C\sqrt x\,\log_2x
 =O(\sqrt x\log x).
\]
最後の項は $o(x)$ なので，
$\psi(x)=\vartheta(x)+o(x)=O(x)$ が従う．
\end{proof}

\section{ゼータ関数と対数微分}

\begin{proposition}[ゼータ関数の有理型延長]
\label{prop:zeta-cont}
$\operatorname{Re}s>1$ で
$\zeta(s)=\sum_{n\ge1}n^{-s}$ と定義する．
\[
 \zeta(s)=\frac{s}{s-1}
 -s\int_1^\infty\{x\}x^{-s-1}\,dx
 \tag{1}
\]
は $\operatorname{Re}s>0$ へ有理型に延長し，
$s=1$ に留数1の単純極を持ち，他の極を持たない．
\end{proposition}
\begin{proof}
$\operatorname{Re}s>1$ で部分和公式を用い
\[
 \zeta(s)=s\int_1^\infty\lfloor x\rfloor x^{-s-1}\,dx.
\]
$\lfloor x\rfloor=x-\{x\}$ と書いて積分すると式 (1) を得る．
$\{x\}$ は有界なので，
式 (1) の積分は $\operatorname{Re}s>0$ の
コンパクト集合上で一様絶対収束する．
従って積分の項はそこで正則であり，
第一項は $s=1$ で留数1の単純極を持つ．
\end{proof}

\begin{lemma}[Euler積の対数微分]\label{lem:logder}
$\operatorname{Re}s>1$ において
\[
 -\frac{\zeta'(s)}{\zeta(s)}
 =\sum_{n=1}^\infty\frac{\Lambda(n)}{n^s}.
 \tag{2}
\]
\end{lemma}
\begin{proof}
Euler積
$\zeta(s)=\prod_p(1-p^{-s})^{-1}$ の対数は
$\log\zeta(s)=\sum_p\sum_{k\ge1}p^{-ks}/k$
である．
$\operatorname{Re}s>1$ のコンパクト集合上で
この級数と導関数級数はともに絶対一様収束する．
従って項別微分でき，
\[
 -\frac{\zeta'(s)}{\zeta(s)}
 =\sum_p\sum_{k\ge1}(\log p)p^{-ks}.
\]
右辺は $\Lambda(n)$ の定義によって式 (2) と一致する．
\end{proof}

\section{境界直線上の零点排除}

\begin{theorem}[Hadamard--de la Vall\'ee-Poussinの零点排除]
\label{thm:no-zero}
任意の実数 $t$ に対し
\[
 \zeta(1+it)\ne0
\]
が成り立つ．ただし $t=0$ では $\zeta$ は極を持ち，
零点ではないという意味である．
\end{theorem}
\begin{proof}
$t\ne0$ を固定する．$\sigma>1$ に対し
Euler積の対数から
\[
 \begin{aligned}
 &\log\!\left(
 \zeta(\sigma)^3
 |\zeta(\sigma+it)|^4
 |\zeta(\sigma+2it)|\right)\\
 &=\sum_p\sum_{k=1}^\infty
 \frac{3+4\cos(kt\log p)+\cos(2kt\log p)}
 {k p^{k\sigma}}.
 \end{aligned}
\]
三角恒等式
$3+4\cos u+\cos2u=2(1+\cos u)^2\ge0$
により右辺は非負である．
従って
\[
 \zeta(\sigma)^3
 |\zeta(\sigma+it)|^4
 |\zeta(\sigma+2it)|\ge1.
 \tag{3}
\]
もし $\zeta(1+it)=0$ なら，
命題\ref{prop:zeta-cont}の正則性によって
$\zeta(\sigma+it)=O(\sigma-1)$ となる．
$\zeta(\sigma+2it)$ は $1+2it\ne1$ で正則なので
$\sigma\downarrow1$ で有界である．
また $\zeta(\sigma)=O((\sigma-1)^{-1})$ である．
従って式 (3) の左辺は
$O((\sigma-1)^{-3+4})=O(\sigma-1)$ となり
零へ向かう．
これは左辺が1以上という不等式と矛盾する．
従って $\zeta(1+it)\ne0$ である．
\end{proof}

\section{NewmanのTauber型補題}

\begin{lemma}[境界正則性からの収束]\label{lem:newman}
$f:[0,\infty)\to\mathbb C$ を有界な局所可積分関数とする．
\[
 F(z)=\int_0^\infty f(t)e^{-zt}\,dt
 \qquad(\operatorname{Re}z>0)
\]
が，閉半平面 $\operatorname{Re}z\ge0$ の
ある開近傍へ正則に延長されると仮定する．
このとき広義積分 $\int_0^\infty f(t)\,dt$ は収束し，
その値は $F(0)$ に等しい．
\end{lemma}
\begin{proof}
$|f(t)|\le B$ とする．
\[
 F_T(z)=\int_0^Tf(t)e^{-zt}\,dt
\]
と置くと $F_T$ は整関数である．
任意に大きい $R>0$ を固定する．
仮定により虚軸のコンパクトな区間
$\{iy:|y|\le R\}$ の近傍で $F$ は正則なので，
十分小さい $\delta>0$ を選び
\[
 D_R=\{z:|z|\le R,\ \operatorname{Re}z\ge-\delta\}
\]
の閉包が $F$ の正則領域に含まれるようにできる．
$\partial D_R$ を正向きの輪郭 $C$ とする．
Cauchyの積分公式から
\[
 F(0)-F_T(0)=\frac1{2\pi i}
 \int_C(F(z)-F_T(z))e^{Tz}
 \left(1+\frac{z^2}{R^2}\right)\frac{dz}{z}.
 \tag{4}
\]
確かに被積分関数の $z=0$ における留数は
$F(0)-F_T(0)$ である．
$C_+=C\cap\{\operatorname{Re}z\ge0\}$，
$C_-=C\cap\{\operatorname{Re}z\le0\}$ と分ける．

$C_+$ 上で $|z|=R$ かつ $\operatorname{Re}z>0$ なら
\[
 |F(z)-F_T(z)|
 =\left|\int_T^\infty f(t)e^{-zt}\,dt\right|
 \le \frac{B e^{-T\operatorname{Re}z}}{\operatorname{Re}z}.
\]
また $|z|=R$ のとき
\[
 \left|\left(1+\frac{z^2}{R^2}\right)\frac1z\right|
 =\frac{2|\operatorname{Re}z|}{R^2}.
 \tag{5}
\]
$C_+$ の長さは高々 $2\pi R$ なので，
式 (4) の $C_+$ 上の積分の絶対値は
高々 $2B/R$ である．

$C_-$ 上の $F_T$ の積分を考える．
$F_T$ は整関数であるから，$C_-$ の左側部分は
半径 $R$ の左半円弧
$C'_-=\{z:|z|=R,\operatorname{Re}z\le0\}$
へ，端点を固定して変形できる．
両輪郭の間の領域には $z=0$ が含まれないので，
この変形で積分値は変わらない．
$C'_-$ 上で $\operatorname{Re}z<0$ なら
\[
 |F_T(z)|
 \le B\int_0^Te^{-t\operatorname{Re}z}\,dt
 \le \frac{B e^{-T\operatorname{Re}z}}{|\operatorname{Re}z|}.
\]
式 (5) と半円弧の長さの評価から，
$C_-$ 上の $F_T$ の寄与も絶対値が
高々 $2B/R$ である．

最後に $C_-$ 上の $F$ の寄与を扱う．
固定した $R,\delta$ に対して $F(z)$ は
$C_-$ 上で有界であり，$z=0$ はこの輪郭上にない．
輪郭のうち端点以外では
$\operatorname{Re}z<0$ なので
$e^{Tz}\to0$ である．
$|e^{Tz}|\le1$ だから，
輪郭を有限個の滑らかな弧に分割して
Lebesgueの収束定理を適用すれば，
$C_-$ 上の $F$ の寄与は $T\to\infty$ で零となる．
従って式 (4) から
\[
 \limsup_{T\to\infty}|F(0)-F_T(0)|
 \le\frac{4B}{R}.
\]
$R>0$ は任意に大きく取れるので，
右辺を零へ近づけられる．
従って $F_T(0)=\int_0^Tf(t)\,dt\to F(0)$ である．
\end{proof}

\section{素数定理の証明}

\begin{lemma}[主要なLaplace変換]\label{lem:laplace}
$t\ge0$ に対して
\[
 f(t)=e^{-t}\psi(e^t)-1
\]
と置く．このとき $f$ は有界であり，
$\operatorname{Re}z>0$ におけるLaplace変換は
\[
 F(z)=-\frac1{z+1}
 \frac{\zeta'(z+1)}{\zeta(z+1)}-\frac1z.
 \tag{6}
\]
さらに $F$ は
$\operatorname{Re}z\ge0$ の開近傍へ正則に延長される．
\end{lemma}
\begin{proof}
補題\ref{lem:chebyshev}によって
$\psi(x)=O(x)$ なので $f$ は有界である．
$\operatorname{Re}z>0$ では
Tonelliの定理または絶対収束する和と積分の交換により
\[
 \begin{aligned}
 \int_0^\infty e^{-(z+1)t}\psi(e^t)\,dt
 &=\sum_{n\ge2}\Lambda(n)
   \int_{\log n}^\infty e^{-(z+1)t}\,dt\\
 &=\frac1{z+1}\sum_{n\ge2}
 \frac{\Lambda(n)}{n^{z+1}}.
 \end{aligned}
\]
補題\ref{lem:logder}と
$\int_0^\infty e^{-zt}\,dt=1/z$ から式 (6) が従う．
$\operatorname{Re}z>0$ ではEuler積から
$\zeta(1+z)\ne0$ である．
境界 $\operatorname{Re}z=0$ で $z\ne0$ なら
定理\ref{thm:no-zero}から
$\zeta(1+z)\ne0$ であり，
$\zeta$ がその点の近傍で正則なことから
式 (6) もそこへ正則に延長できる．
$z=0$ では
$\zeta(1+z)=z^{-1}+a_0+a_1z+\cdots$ と展開されるので
\[
 -\frac{\zeta'(1+z)}{\zeta(1+z)}
 =\frac1z+\text{正則関数}
\]
となる．
これを式 (6) に代入すると
$1/(z(1+z))-1/z=-1/(1+z)$ となり，
極が相殺されるため $F$ は $z=0$ でも正則である．
従って閉半平面の各点の近傍で正則に延長され，
それらを合わせた開近傍で正則となる．
\end{proof}

\begin{theorem}[Chebyshev関数の漸近式]\label{thm:psi}
\[
 \psi(x)\sim x\qquad(x\to\infty).
\]
\end{theorem}
\begin{proof}
補題\ref{lem:laplace}の $f$ は
補題\ref{lem:newman}の仮定を満たすので
$\int_0^\infty f(t)\,dt$ は収束する．
$x=e^t$ と変数変換すると
\[
 \int_1^\infty\frac{\psi(x)-x}{x^2}\,dx
 \tag{7}
\]
が広義積分として収束することに等しい．

もしある $\varepsilon>0$ について，
$\psi(X)\ge(1+\varepsilon)X$ となる
任意に大きい $X$ が存在するとする．
$1<\lambda<1+\varepsilon$ を固定し，
$\psi$ は単調非減少だから
$X\le x\le\lambda X$ で
$\psi(x)-x\ge(1+\varepsilon)X-x\ge0$ である．
従って
\[
 \int_X^{\lambda X}\frac{\psi(x)-x}{x^2}\,dx
 \ge\int_1^\lambda\frac{1+\varepsilon-u}{u^2}\,du>0
\]
となり，右辺は $X$ に依存しない．
しかし式 (7) の収束から，
$X\to\infty$ では左辺は零へ向かうはずである．
矛盾する．
同様に，ある $\varepsilon>0$ について
$\psi(X)\le(1-\varepsilon)X$ が
任意に大きい $X$ で成立すると仮定する．
$1-\varepsilon<\lambda<1$ を取ると
$\lambda X\le x\le X$ で
$\psi(x)-x\le(1-\varepsilon)X-x<0$ である．
従って
\[
 \int_{\lambda X}^{X}
 \frac{\psi(x)-x}{x^2}\,dx
 \le\int_\lambda^1
 \frac{1-\varepsilon-u}{u^2}\,du<0
\]
となり，式 (7) の収束に反する．
以上から $\psi(x)/x\to1$ が従う．
\end{proof}

\begin{theorem}[素数定理]\label{thm:pnt}
\[
 \boxed{\displaystyle
 \pi(x)\sim\frac{x}{\log x}}
 \qquad(x\to\infty)
\]
が成り立つ．
\end{theorem}
\begin{proof}
補題\ref{lem:prime-powers}と
補題\ref{lem:chebyshev}によって
$\psi(x)-\vartheta(x)=o(x)$ だから，
定理\ref{thm:psi}から
$\vartheta(x)\sim x$ を得る．
Abelの部分和公式により
\[
 \pi(x)=\sum_{p\le x}\frac{\log p}{\log p}
 =\frac{\vartheta(x)}{\log x}
  +\int_2^x\frac{\vartheta(t)}
                 {t(\log t)^2}\,dt.
 \tag{8}
\]
第一項は $\sim x/\log x$ である．
第二項は補題\ref{lem:chebyshev}から
\[
 O\left(\int_2^x\frac{dt}{(\log t)^2}\right)
 =O\left(\frac{x}{(\log x)^2}\right)
 =o\left(\frac{x}{\log x}\right)
\]
である．
真ん中の評価は積分区間を
$[2,\sqrt x]$ と $[\sqrt x,x]$ に分けて確かめられる：
前者は $O(\sqrt x)$，後者は
$O(x/(\log x)^2)$ である．
従って式 (8) の第一項だけが主項となり，
求める漸近式を得る．
\end{proof}

\begin{remark}
本証明ではゼータ関数の関数等式や
臨界帯における零点分布を用いない．
必要となる解析的事実は，
$s=1$ の極，$\operatorname{Re}s=1$ 上の零点の不在，
および境界正則性を積分の収束へ変換する
Tauber型補題である．
\end{remark}

\section{到達課題の完全解答}
\noindent\textbf{素数計数関数から重み付き総和へ．}
\begin{proposition}
$\pi(x)\sim x/\log x$ ならば
\[
\vartheta(x)=\sum_{p\le x}\log p\sim x
\]
である．
\end{proposition}
\begin{proof}
Abelの部分和公式より
\[
\vartheta(x)
=\pi(x)\log x-\int_2^x\frac{\pi(t)}t\,dt.
\]
第一項は $x(1+o(1))$ である．
第二項について $\pi(t)\le Ct/\log t$ が十分大きな
$t$ で成り立つから
\[
\int_2^x\frac{\pi(t)}t\,dt
=O\left(\int_2^x\frac{dt}{\log t}\right)
=O(x/\log x)=o(x).
\]
最後の評価は積分区間を
$[2,\sqrt x]$ と $[\sqrt x,x]$ に分ければよい．
\end{proof}
\noindent\textbf{到達確認．}
$\pi(x)\sim x/\log x$ から
$\sum_{p\le x}1/p$ の発散は導けるか．
\par\noindent\textbf{解答．}
部分和公式から
\[
\sum_{p\le x}\frac1p
=\frac{\pi(x)}x+\int_2^x\frac{\pi(t)}{t^2}\,dt.
\]
ある $T$ 以降で $\pi(t)\ge t/(2\log t)$ なので，
右辺の積分は
$\frac12\int_T^x dt/(t\log t)$ 以上となり発散する．
\end{document}